%%%%%
%
%% Dans ce fichier, les �ventuels accents sont cod�s
%% en ISO-8859-15 (aka latin9)
%% La ligne suivante donne trois exemples 
%% e accent aigu : �, e accent grave : �, e-dans-l'o : �
\section{Construction d'un vecteur "tableau" \`a partir de 30 vecteurs scalaires}

\itemstyle{Op\'erateur~:} \operfmt{scalar30\_to\_tab}
\index{scalar30\_to\_tab}

%%%%%OBJET%%%%
\itemstyle{Objet~:}
Construction d'un vecteur "tableau" \`a partir de 30 vecteurs scalaires
de type \texttt{float}

%%%%%ARGUMENTS%%%%
\bigskip
%% Dans ce fichier, les �ventuels accents sont cod�s
%% en ISO-8859-15 (aka latin9)
%% La ligne suivante donne trois exemples 
%% e accent aigu : �, e accent grave : �, e-dans-l'o : �
%%%%%ENTREES%%%%
\begin{tabarg}
\multicolumn{3}{l}{\titretab {Entr\'ees}}\\
\hline
$X0$ & Vecteur, \texttt{float} & premier param�tre "scalaire"\\
$X1$ & Vecteur, \texttt{float} & second param�tre "scalaire"\\
\vdots & \vdots & \vdots \\
$X{29}$ & Vecteur, \texttt{float} & trenti�me param�tre scalaire \\
\end{tabarg}

\medskip
%%%%%SORTIES%%%%
\begin{tabarg}
\multicolumn{3}{l}{\titretab {Sorties}}\\
\hline
$Z$ & Vecteur, \texttt{float[30]} & vecteur "1D" de taille 30 sur la deuxi�me dimension\\
\end{tabarg}


%%%%%FORMULE%%%%
\itemstyle{Formule}

%% Formule de l'operateur : scalar30_to_tab
%% Dans ce fichier, les �ventuels accents sont cod�s
%% en ISO-8859-15 (aka latin9)
%% La ligne suivante donne trois exemples 
%% e accent aigu : �, e accent grave : �, e-dans-l'o : �
\[
Z_k = Xk  \qquad k=0,29
\]



%%%%%COMMENTAIRE%%%%
\itemstyle{Commentaire~:}
Cet op�rateur ne renseigne pas les attibuts "bornes" et "unit\_bornes"
du vecteur cr��.

%%%%%AUTEUR%%%%
\itemstyle{Auteur~:}
CNRM/GMEI/TRAMM

%%%%%REFERENCE%%%%
\itemstyle{R\'ef\'erences~:}
